Agents that share context can talk privately in public: an overhearer without that context cannot decode what they mean.
We study whether a small language model can be trained to do this. Agent A knows a two-part secret (who will bring the birthday person, and where the party is)
and must convey it to B, with whom it shares eight private memories, while chatting about an unrelated topic in front of C, who lacks the memories.
Memories are generated fresh for every episode before the secret is drawn, so no reusable code exists.
We warm-start a 3B speaker by supervised fine-tuning on filtered conversations from a 14B model and then train it with multi-turn GRPO.
The reward favours conversations in which a memory-aware listener is confident in the truth and a memory-blind listener is not, provided the chat is coherent and names nothing.
Our original design trained 1.5B listeners and retrained the outsider every round; the trained listeners did not exceed a guessing baseline, so both listeners are the same 14B model prompted zero-shot, with a held-out outsider prompt for evaluation. We report the original plan, the failed attempts and the revisions.
On \x{nTest} held-out test episodes, the warm-start speaker gives the insider an advantage of \res{S1}{insider_adv} points in both-correct accuracy, against \res{S0}{insider_adv} for the untrained speaker, whose turns leak names \res{S0}{leak_rate}\% of the time.
GRPO raised the coherence pass rate but did not improve the insider advantage (\res{S3}{insider_adv}); the gain is concentrated in the helper slot, and a held-out outsider prompt confirms the outsider's low accuracy.
